\documentclass[11pt,a4paper]{article}
\usepackage[margin=2.4cm]{geometry}
\usepackage{amsmath,amssymb,amsthm,mathtools}
\usepackage{booktabs}
\usepackage{enumitem}
\usepackage{xcolor}
\usepackage{hyperref}

\hypersetup{
  colorlinks=true,
  linkcolor=blue!70!black,
  citecolor=blue!70!black,
  urlcolor=blue!70!black,
  pdftitle={Cuboctahedral Rotator Calculus: Discrete Euler Exponential on Z_512, Constant Wronskian, and Closed-Form Quadrature on the Cuboctahedron O_2(3) in Lean 4},
  pdfauthor={Dmitrii A. Karataev}
}

\theoremstyle{plain}
\newtheorem{theorem}{Theorem}[section]
\newtheorem{lemma}[theorem]{Lemma}
\newtheorem{proposition}[theorem]{Proposition}
\newtheorem{corollary}[theorem]{Corollary}
\theoremstyle{definition}
\newtheorem{definition}[theorem]{Definition}
\newtheorem{remark}[theorem]{Remark}

\title{\textbf{Cuboctahedral Rotator Calculus:\\ Discrete Euler Exponential on $\mathbb{Z}_{512}$, Constant Wronskian, and Closed-Form Quadrature on the Cuboctahedron $O_2(3)$ in Lean~4}}
\author{Dmitrii A.~Karataev\thanks{AI-Guiders Research (\url{https://github.com/AI-Guiders}). Formal Lean~4 verification and analytical development were conducted in partnership with \emph{Luch} (AI research partner, AI-Guiders).}}
\date{}

\begin{document}
\maketitle

\begin{abstract}
We construct an exact, pole-free $C^\infty$ projective differential, variational, and integral calculus over the fields $\mathbb{Q}(i)$ and $\mathbb{Q}(\sqrt{2})$ parameterized by a cyclic residue register $\mathbb{Z}_{4Q}$ ($N_{\mathrm{ring}} = 4\,Q_\square = 512$ at $Q_\square = 2^{|\triangle|-1} = 128$ on the three-dimensional cuboctahedral shell $O_2(3)\subset\{-1,0,+1\}^3$ with $|V_{O_2}|=12$ vertices, $|\square|=6$ square faces, and $|\triangle|=8$ triangular faces). In classical analysis, the Weierstrass substitution $t = \tan(\theta/2) = S/C$ suffers from a pole $t\to\infty$ at $\theta=\pi$ ($C=0$), while in $L^1$ ``taxicab trigonometry'' (Krause 1986; Thompson, \texttt{arXiv:1205.3814}) piecewise derivatives break down at the four diamond corners and naive modulo-$4Q$ residue multiplication has zero divisors. We resolve both obstacles simultaneously by keeping the rotator pair $\mathcal{E}_Q(\varphi) = (\mathcal{C}_Q(\varphi),\mathcal{S}_Q(\varphi))$ homogeneous on the closed $L^1$-diamond $|\mathcal{C}_Q(\varphi)|+|\mathcal{S}_Q(\varphi)|\equiv Q\ge 1$ and lifting it via $\pi_{M\to T}(C,S) = (C^2-S^2 : 2CS : C^2+S^2)$ into projective coordinates over $\mathbb{Q}(i)$ and $\mathbb{Q}(\sqrt{2})$. Because $2(\mathcal{C}_Q^2+\mathcal{S}_Q^2) = (\mathcal{C}_Q+\mathcal{S}_Q)^2 + (\mathcal{C}_Q-\mathcal{S}_Q)^2 \ge Q^2 \ge 1$, the Weierstrass denominator $R_Q(\varphi) = \mathcal{C}_Q(\varphi)^2+\mathcal{S}_Q(\varphi)^2 \ge 1$ never vanishes (\texttt{thm34}), eliminating zero divisors and the $\theta=\pi$ pole, while at the four diamond corners the derivative jump multiplies a vanishing coordinate so that the one-step difference Wronskian $\mathcal{W}_Q(\varphi)=\mathcal{C}_Q(\varphi)\Delta\mathcal{S}_Q(\varphi)-\mathcal{S}_Q(\varphi)\Delta\mathcal{C}_Q(\varphi)\equiv Q$ is identically constant across all $4Q$ points (\texttt{thm1}).

From this projective field structure, all differential, variational, and integral laws follow directly: (1)~the difference derivative $(Df)(x)=f(x+1)-f(x)$ satisfies the projective quotient-derivative Wronskian law $f(x+1)g(x)-f(x)g(x+1) = -\mathcal{W}(f,g)(x)$, the variational Euler--Lagrange derivative $\delta\mathcal{S}/\delta D_w = P_1 Q_2$ on trajectory monoids (\texttt{thm35}), and the inductive $C^\infty$ higher-derivative tower $D^n f$ (\texttt{thm38}) with pole-free denominators $\prod_{j=0}^n R_Q(\varphi+j)\ge 1$ for all $n\in\mathbb{N}$ and exact remainder-free Taylor--Newton series; (2)~the two-tick symplectic Verlet recurrence and two-ring coupling at the cuboctahedron modulus $k^2 = 1/\|e\|_2^2 = 1/2$ yield closed-form integer evaluations of the Fresnel diffraction integrals ($2048$), lemniscate Jacobi elliptic integrals, the Dirichlet sine integral ($128$), Parseval energy ($2,796,288$), and the zero-order Bessel integral ($\mathcal{J}_{0,3}(0)=65,536$, $\mathcal{J}_{0,3}(4)=0$, \texttt{thm2}); and (3)~slicing the $12$ vertices of $O_2(3)$ along any axis yields the equizonal partition $(4,4,4)$, which with the origin weight $12$ forms Simpson's weights $(4,16,4)=4\times(1,4,1)$ and reproduces the exact 3D continuum Riemann integral $\iiint_{[-1,1]^3} F\,\mathrm{d}V$ in $13$ points (\texttt{thm3}). All theorems (\texttt{thm1}, \texttt{thm2}, \texttt{thm3}, \texttt{thm34}, \texttt{thm35}, \texttt{thm38}) are formally verified in \textbf{Lean~4} (\texttt{0 sorryAx}, \texttt{0 Classical.choice}).

\medskip\noindent
\textbf{Keywords:} piecewise-linear trigonometry, Euler's identity, difference Wronskian, Weierstrass substitution, cuboctahedron, Fresnel integrals, lemniscate Jacobi elliptic functions, formal verification, Lean~4.

\smallskip\noindent
\textbf{MSC 2020:} 26A06, 33E05, 39A12, 65D30, 68V20.
\end{abstract}

\section{Introduction and Historical Context}
\label{sec:intro}

The components of the calculus developed in this paper have deep roots in four classical mathematical traditions that evolved independently and were never synthesized into a single computational framework:
\begin{enumerate}[label=\textbf{(\arabic*)},leftmargin=*,noitemsep]
\item \textbf{Projective parameterization of genus-0 curves vs.\ the affine Weierstrass pole (Diophantus, Brahmagupta 628, Pl\"ucker, Weierstrass):}
  in standard analysis, the universal trigonometric substitution $t = \tan(\theta/2) = S/C \mapsto \bigl((1-t^2)/(1+t^2),\,2t/(1+t^2)\bigr)$ is written in a single affine coordinate $t\in\mathbb{Q}\cup\{\infty\}$ and therefore suffers from a pole $t\to\infty$ at $\theta=\pi$ ($C=0$), requiring two coordinate charts. In projective algebraic geometry, by contrast, the homogeneous map $\mathbb{P}^1(\mathbb{Q})\to S^1(\mathbb{Q})$, $(C : S)\mapsto (C^2-S^2 : 2CS : C^2+S^2)$, is regular everywhere because $C^2+S^2$ is anisotropic over $\mathbb{Q}$. Yet projective equivalence classes $(C : S)\sim(\lambda C : \lambda S)$ lacked a canonical finite cyclic clock section with a uniform step $\varphi\mapsto\varphi+1$;
\item \textbf{Exact calculus of finite differences and remainder-free Taylor--Newton series (Gregory 1670, Newton 1687, Taylor 1715, Boole 1860~\cite{Boole1860}):}
  the difference derivative $(Df)(x)=f(x+1)-f(x)$, the quotient rule $D(f/g) = -\mathcal{W}(f,g)/(g\,g_{+1})$, Pochhammer falling factorials $x^{\underline{n}}$, and exact finite Taylor--Newton expansions were fully established by Newton, Taylor (\emph{Methodus Incrementorum}, 1715), and Boole (\emph{A Treatise on the Calculus of Finite Differences}, 1860~\cite{Boole1860}). In the 20th century, however, scientific computing abandoned exact projective difference algebra in favor of IEEE~754 floating-point approximations (`Float64`) with truncation and rounding drift;
\item \textbf{Piecewise-linear $L^1$ ``taxicab trigonometry'' (Minkowski, Krause 1986~\cite{Krause1986}, Thompson 2011--2012~\cite{Thompson2011,Thompson2012}):}
  parameterizing the $L^1$-diamond $|C|+|S|=Q$ by arc length provided a natural cyclic phase register $\varphi\in\mathbb{Z}_{4Q}$, but in~\cite{Thompson2012} (\texttt{arXiv:1205.3814}) three obstacles remained unresolved: (a)~piecewise derivatives broke down at the four diamond corners and lost Euclidean coupling; (b)~quadrant transitions required branched powers $i^{k-2}$; and (c)~integral calculus, the lift to the Euclidean circle, and classical non-elementary integrals (Fresnel, Jacobi elliptic, Bessel, Dirichlet) were left open;
\item \textbf{Invariant 3D cubature on the cuboctahedron $O_2(3)$ (Stroud 1971~\cite{Stroud1971}, Sobolev 1974~\cite{Sobolev1974}) and Mal'cev's algebraic systems (1970~\cite{Malcev1973}):}
  spherical and polyhedral cubature theory~\cite{Sobolev1974,Stroud1971} established isotropic 3D integration on the $12$ vertices of the cuboctahedron plus the origin, but treated the weights in continuum floating-point arithmetic rather than linking $O_2(3)$ to the quadratic field $\mathbb{Q}(\sqrt{2})$ and the $\mathbb{Z}_{4Q}$ rotator carrier in the model-theoretic sense of Mal'cev~\cite{Malcev1973}.
\end{enumerate}

In this paper (Part~I of the mathematical series), we unite these four traditions by choosing the closed $L^1$-diamond $|C(\varphi)|+|S(\varphi)|=Q\ge 1$ on $\mathbb{Z}_{4Q}$ as the canonical global gauge section of the projective line $\mathbb{P}^1(\mathbb{Q})$ over the three-dimensional cuboctahedral shell $O_2(d) = \{e=(e_x,e_y,e_z)\in\{-1,0,+1\}^3 \mid e_x^2+e_y^2+e_z^2 = d-1\}$ ($d=3$) and phase ring $\mathbb{Z}_{N_{\mathrm{ring}}(d)}$:
\[
  |{\triangle}(d)| = 2^d = 8,\qquad
  |\square(d)| = 2d = 6,\qquad
  |V_{O_2}(d)| = 2d(d-1) = 12,
\]
\[
  d_{\mathrm{in}}(\square) = 4,\qquad
  d_{\mathrm{out}}(E)(d) = 2^d - d = 5,\qquad
  \|e\|_2^2 = d - 1 = 2,
\]
\[
  Q_\square(d) = 2^{|{\triangle}(d)|-1} = 128,\qquad
  \Phi_{\mathrm{half}}(d) = 2\,Q_\square(d) = 256,\qquad
  N_{\mathrm{ring}}(d) = 4\,Q_\square(d) = 512,\qquad
  K_P(d) = \frac{N_{\mathrm{ring}}(d)}{|{\triangle}(d)|} = 64.
\]
All definitions and theorems below correspond to \texttt{thm1}, \texttt{thm2}, \texttt{thm3}, \texttt{thm34}, \texttt{thm35}, and \texttt{thm38} in the self-contained \textbf{Lean~4}~\cite{Moura2021} module \texttt{RotatorQuadrature.lean}.

\begin{definition}[Model-Theoretic Two-Sorted Algebraic System $\langle\mathcal{M},\mathcal{T}\rangle$ in the Sense of Mal'cev~\cite{Malcev1973}]
\label{def:malcev-carriers}
Throughout this series, the term \emph{carrier} (\emph{носитель} in the Russian model-theoretic tradition of A.\,I.~Mal'cev, \emph{Algebraic Systems}~\cite{Malcev1973}) is used strictly in its standard model-theoretic meaning: the underlying domain set of an algebraic structure $\mathfrak{A} = \langle A;\,\Sigma\rangle$. We formulate the calculus as a two-sorted algebraic structure $\mathfrak{S} = \langle \mathcal{M},\,\mathcal{T};\,\pi_{M\to T},\,\iota_{T\to M}\rangle$ with two carriers:
\begin{enumerate}[leftmargin=*,noitemsep]
\item the \textbf{discrete $M$-carrier} $\mathcal{M} = \mathbb{Z}_{512}$ (equipped with $L^1$ diamond projections $\mathcal{C}_3,\mathcal{S}_3$ and modular arithmetic over the cuboctahedral lattice $\mathbb{Z}^3$);
\item the \textbf{continuous (Euclidean) $T$-carrier} $\mathcal{T}$ (the continuum sort of $L^2$ Euclidean circle invariants $X^2+Y^2=R^2$, rational/real projective coordinates $\mathbb{Q}\subset\mathbb{R}$, and continuum Riemann integrals);
\item connected by the root-free Weierstrass lift $\pi_{M\to T}(C,S) = (C^2-S^2,2CS,C^2+S^2)$, linear retraction $\iota_{T\to M}$, and stereographic homeomorphism $[0,512]/512 = [0,1] \cong \mathbb{R}$.
\end{enumerate}
\end{definition}

\section{Discrete Euler Exponential, Euler's Identity \texorpdfstring{$e^{i\pi}+1=0$}{e\^{}(i pi)+1=0}, Constant Wronskian, and the Weierstrass 512-gon Lift}
\label{sec:wronskian-euler}

For any $x\in\mathbb{Z}$, define the principal symmetric branch wrapper into $(-\Phi_{\mathrm{half}}(d),\Phi_{\mathrm{half}}(d)]$:
\begin{equation}
  \label{eq:wrap-ring}
  \mathrm{wrap}_d(x) \;:=\; \Phi_{\mathrm{half}}(d) - \bigl((\Phi_{\mathrm{half}}(d) - x) \bmod N_{\mathrm{ring}}(d)\bigr),
\end{equation}
and the branchless piecewise-linear cosine and sine projections:
\begin{equation}
  \label{eq:proj-cs}
  \mathcal{C}_d(\varphi) \;:=\; Q_\square(d) - |\mathrm{wrap}_d(\varphi)|,\qquad
  \mathcal{S}_d(\varphi) \;:=\; \mathcal{C}_d\bigl(\varphi - Q_\square(d)\bigr).
\end{equation}

\begin{definition}[Discrete Euler exponential and one-step difference Wronskian]
\label{def:euler-exp-wronskian}
We define the \emph{discrete Euler exponential} on $\mathbb{Z}_{N_{\mathrm{ring}}(d)}$ by the pair
\[
  \mathcal{E}_d(\varphi) \;:=\; \bigl(\mathcal{C}_d(\varphi),\,\mathcal{S}_d(\varphi)\bigr) \;\in\; \mathbb{Z}^2,
\]
and the \emph{one-step difference Wronskian} by
\begin{equation}
  \label{eq:rotator-wronskian}
  \mathcal{W}_d(\varphi) \;:=\; \mathcal{C}_d(\varphi)\,\mathcal{S}_d(\varphi+1) - \mathcal{S}_d(\varphi)\,\mathcal{C}_d(\varphi+1)
  \;=\; \mathcal{C}_d(\varphi)\,\Delta\mathcal{S}_d(\varphi) - \mathcal{S}_d(\varphi)\,\Delta\mathcal{C}_d(\varphi),
\end{equation}
where $\Delta f(\varphi) = f(\varphi+1)-f(\varphi)$.
\end{definition}

To each $\varphi\in\mathbb{Z}_{N_{\mathrm{ring}}(d)}$ we associate the rational Brahmagupta--Weierstrass lift onto the Euclidean circle $X^2+Y^2=R^2$:
\begin{equation}
  \label{eq:weierstrass-lift}
  X_d(\varphi) := \mathcal{C}_d(\varphi)^2 - \mathcal{S}_d(\varphi)^2,\qquad
  Y_d(\varphi) := \|e\|_2^2\,\mathcal{C}_d(\varphi)\,\mathcal{S}_d(\varphi),\qquad
  R_d(\varphi) := \mathcal{C}_d(\varphi)^2 + \mathcal{S}_d(\varphi)^2.
\end{equation}

\begin{theorem}[Discrete Euler's Identity, Multi-Turn Periodicity, and Constant Difference Wronskian]
\label{thm:main-1}
For any quarter-period $Q\in\mathbb{Z}$, intra-quadrant phase $m\in\mathbb{Z}$, turn count $k_{\mathrm{turns}}\in\mathbb{Z}$, and integration horizon $L\in\mathbb{N}$, the piecewise-linear rotator $(\mathcal{C}_Q,\mathcal{S}_Q)$ on $\mathbb{Z}_{4Q}$ satisfies (\texttt{thm1\_constant\_wronskian\_weierstrass\_lift\_and\_newton\_leibniz}):
\begin{enumerate}[label=\textbf{(\roman*)},leftmargin=*,noitemsep]
\item \textbf{Multi-turn periodicity and Euler's identity $e^{i\pi}+1=0$:}
  $\mathcal{E}_d(\varphi + 4Qk_{\mathrm{turns}}) = \mathcal{E}_d(\varphi)$, the $L^1$ norm is identically constant $|\mathcal{C}_Q(\varphi)| + |\mathcal{S}_Q(\varphi)| \equiv Q$, the quarter-turn shift $+Q$ acts as exact multiplication by the imaginary unit~$i$: $\mathcal{E}_Q(\varphi + Q) = (-\mathcal{S}_Q(\varphi),\,+\mathcal{C}_Q(\varphi))$, and the half-turn shift $+2Q$ satisfies $\mathcal{E}_Q(\varphi + 2Q) + \mathcal{E}_Q(\varphi) = (0,0)$ for all $\varphi\in\mathbb{Z}$.
\item \textbf{Constant difference Wronskian and pointwise-constant integral extraction:}
  across all four quadrants and all four corner vertices of the $L^1$ diamond, the one-step difference Wronskian is identically equal to the constant $Q$:
  \begin{equation}
    \label{eq:constant-wronskian-all-512}
    \mathcal{W}_Q(\varphi) \;=\; (Q-m)(m+1) - m(Q-m-1) \;\equiv\; Q
    \qquad (\forall Q,m\in\mathbb{Z}).
  \end{equation}
  Consequently, for any pointwise-constant invariant $g(\varphi)\equiv c_{\mathrm{inv}}$ (including $|\mathcal{C}_Q(\varphi)|+|\mathcal{S}_Q(\varphi)|\equiv Q$ and $\mathcal{W}_Q(\varphi)\equiv Q$), the constant factors directly out of the discrete integral over any horizon $L\in\mathbb{N}$:
  \[
    \sum_{\varphi=0}^{L-1} g(\varphi) \;=\; c_{\mathrm{inv}}\sum_{\varphi=0}^{L-1} 1 \;=\; L\cdot c_{\mathrm{inv}},
    \qquad
    \sum_{\varphi=0}^{L-1} c_{\mathrm{sq}}\,f(\varphi) \;=\; c_{\mathrm{sq}}\sum_{\varphi=0}^{L-1} f(\varphi).
  \]
\end{enumerate}
\end{theorem}

\begin{theorem}[Root-Free Euclidean Weierstrass Lift and Brahmagupta Step Identities]
\label{thm:main-1b}
For all integer rotator projections $\mathcal{C}_0,\mathcal{S}_0,\mathcal{C}_1,\mathcal{S}_1\in\mathbb{Z}$, the rational Weierstrass coordinates~\eqref{eq:weierstrass-lift} satisfy the exact Euclidean circle identity $X^2 + Y^2 = (\mathcal{C}_0^2-\mathcal{S}_0^2)^2 + (2\mathcal{C}_0\mathcal{S}_0)^2 = (\mathcal{C}_0^2+\mathcal{S}_0^2)^2 = R^2$ and the Brahmagupta--Weierstrass step decomposition with inner product $\mathcal{I} = \mathcal{C}_0\mathcal{C}_1 + \mathcal{S}_0\mathcal{S}_1$ and Wronskian $\mathcal{W} = \mathcal{C}_0\mathcal{S}_1 - \mathcal{S}_0\mathcal{C}_1$:
\begin{equation}
  \label{eq:brahmagupta-step}
  R_0 R_1 - \mathcal{I}^2 \;=\; \mathcal{W}^2,\qquad
  X_0 Y_1 - Y_0 X_1 \;=\; 2\,\mathcal{W}\,\mathcal{I},\qquad
  X_0 X_1 + Y_0 Y_1 \;=\; \mathcal{I}^2 - \mathcal{W}^2.
\end{equation}
\end{theorem}

\begin{theorem}[Fundamental Theorem of Discrete Calculus and Multiplierless Arithmetic]
\label{thm:main-1c}
For every dimension $d\in\mathbb{N}$, upper limit $L\in\mathbb{N}$, turn count $k_{\mathrm{turns}}\in\mathbb{Z}$, function $f\colon\mathbb{Z}\to\mathbb{Z}$, and integers $Q,x,y\in\mathbb{Z}$ (\texttt{thm1\_constant\_wronskian\_weierstrass\_lift\_and\_newton\_leibniz}):
\begin{equation}
  \label{eq:newton-leibniz-universal}
  \sum_{\varphi=0}^{L} f(\varphi) - \sum_{\varphi=0}^{L-1} f(\varphi) \;=\; f(L),\qquad
  \Delta\widehat{\Phi}_{\mathcal{C}}(d,L) \;=\; \mathcal{C}_d(L),\qquad
  \Delta\widehat{\Phi}_{\mathcal{S}}(d,L) \;=\; \mathcal{S}_d(L),
\end{equation}
with full-turn vanishing $k_{\mathrm{turns}}\widehat{\Phi}_{\mathcal{C}}(3,512) = k_{\mathrm{turns}}\widehat{\Phi}_{\mathcal{S}}(3,512) = 0$, multi-turn half-wave area $k_{\mathrm{turns}}\widehat{\Phi}_{\mathcal{S}}(3,256) = k_{\mathrm{turns}}Q_\square(3)^2 = 16,384\,k_{\mathrm{turns}}$, and quarter-square multiplierless multiplication $(x+y)^2 - (x-y)^2 = 4xy$.
\end{theorem}

\begin{proof}[Proof of Theorems~\ref{thm:main-1}--\ref{thm:main-1c}]
In~\eqref{eq:constant-wronskian-all-512}, expanding $(Q-m)(m+1) - m(Q-m-1) = (Qm + Q - m^2 - m) - (Qm - m^2 - m) = Q$ cancels all quadratic and linear terms in $m$ for arbitrary $Q,m\in\mathbb{Z}$; at the diamond corner $m=Q$, $0\cdot(Q-1) - Q\cdot(-1) = Q$. Inducting on $L\in\mathbb{N}$ gives $\sum_{\varphi=0}^{L-1} g(\varphi) = L\,c_{\mathrm{inv}}$ whenever $g(\varphi)\equiv c_{\mathrm{inv}}$. The polynomial identities of Theorems~\ref{thm:main-1b}--\ref{thm:main-1c} hold over $\mathbb{Z}$ (\texttt{omega} in Lean~4).
\end{proof}

\begin{theorem}[Projective Field Structure of the Rotator Carrier: Absence of Zero Divisors on the $L^1$-Diamond, Field Invertibility in $\mathbb{Q}(i)$ and $\mathbb{Q}(\sqrt{2})$, and the $C^\infty$ Derivative and Variational Tower, \texttt{thm34}, \texttt{thm35}, \texttt{thm38}]
\label{thm:main-34-field}
Unlike naive residue multiplication modulo composite $4Q$ (which possesses zero divisors), the rotator pair $(C,S) = (\mathcal{C}_Q(\varphi),\mathcal{S}_Q(\varphi))$ on the $L^1$-diamond $|C|+|S|=Q\ge 1$ and its Weierstrass lift $(X_T,Y_T,R_T) = (C^2-S^2,2CS,C^2+S^2)$ operate in homogeneous projective coordinates over the fields $\mathbb{Q}(i)$ and $\mathbb{Q}(\sqrt{2})$ (\texttt{thm34\_rotator\_projective\_field\_invertibility\_zero\_divisor\_freedom\_and\_q\_i\_sqrt2\_fields}):
\begin{enumerate}[label=\textbf{(\roman*)},leftmargin=*,noitemsep]
\item \textbf{Strict norm positivity and absence of zero divisors:}
  $2(C^2+S^2) = (C+S)^2 + (C-S)^2 \ge Q^2 \ge 1$, so $R_T(\varphi) = C^2+S^2 \ge 1 > 0$ at every phase $\varphi\in\mathbb{Z}_{4Q}$; by Brahmagupta--Fibonacci norm multiplicativity $(C_1 C_2 - S_1 S_2)^2 + (C_1 S_2 + S_1 C_2)^2 = (C_1^2+S_1^2)(C_2^2+S_2^2) \ge 1$, the product of any two rotator states is never zero;
\item \textbf{Abelian additive group, Abelian multiplicative group of non-zero elements, and distributivity in $\mathbb{Q}(i)$:}
  in homogeneous triples $(X,Y,R)\in\mathbb{Z}^3$ representing $(X+iY)/R$, \texttt{gaussFieldAdd} and \texttt{gaussFieldMul} satisfy commutativity, associativity, and distributivity, $(X,Y,R)\oplus(-X,-Y,R) = (0,0,R^2)$, and the multiplicative inverse $\texttt{gaussFieldInv}(X,Y,R) = (RX,-RY,X^2+Y^2)$ satisfies
  \begin{equation}
    \label{eq:gauss-field-inverse}
    (X,Y,R)\otimes (RX,\,-RY,\,X^2+Y^2) \;=\; \bigl(R(X^2+Y^2),\,0,\,R(X^2+Y^2)\bigr) \;\sim\; (1,0,1);
  \end{equation}
\item \textbf{Field axioms, Pell norm multiplicativity, and multiplicative inverse in $\mathbb{Q}(\sqrt{2})$:}
  in homogeneous triples $(a,b,D)\in\mathbb{Z}^3$ representing $(a+b\sqrt{2})/D$, \texttt{sqrt2FieldAdd} and \texttt{sqrt2FieldMul} satisfy commutativity and associativity, $(a_1 a_2 + 2b_1 b_2)^2 - 2(a_1 b_2 + b_1 a_2)^2 = (a_1^2-2b_1^2)(a_2^2-2b_2^2)$, and $\texttt{sqrt2FieldInv}(a,b,D) = (Da,-Db,a^2-2b^2)$ satisfies $(a,b,D)\otimes_{\sqrt{2}}(Da,-Db,a^2-2b^2) = \bigl(D(a^2-2b^2),\,0,\,D(a^2-2b^2)\bigr)\sim(1,0,1)$;
\item \textbf{General field derivative operator $Df$, quotient-derivative Wronskian law, variational derivative $\delta\mathcal{S}/\delta D_w = P_1 Q_2$, and $C^\infty$ Taylor--Newton tower (\texttt{thm35}, \texttt{thm38}):}
  for any $f,g\colon\mathbb{Z}\to\mathbb{Z}$, the difference derivative $(Df)(x)=f(x+1)-f(x)$ satisfies $D(af+bg)=a\,Df+b\,Dg$, the Leibniz rule $D(f\cdot g)(x)=f(x+1)Dg(x)+Df(x)g(x)$, $D(Df)(x)=f(x+2)-2f(x+1)+f(x)$, and the projective quotient-derivative law $f(x+1)g(x)-f(x)g(x+1) = g(x)Df(x)-f(x)Dg(x) = -\mathcal{W}(f,g)(x)$ (so $f/g$ is stationary at $x \iff \mathcal{W}(f,g)(x)=0$ and strictly increases $\iff \mathcal{W}(f,g)(x)<0$), while on any composite trajectory $(L_1,w,L_2)$ with action $\mathcal{S} = Q_2(Q_w D_1 + P_1 D_w) + P_1 P_w D_2$, varying $D_w\mapsto D_w'$ at fixed boundary multipliers $(P_w,Q_w)$ yields $\delta\mathcal{S} = P_1(D_w'-D_w)Q_2$ with Euler--Lagrange stationarity $\delta\mathcal{S}=0 \iff D_w'=D_w$ for $P_1,Q_2\ge 1$. Moreover, the inductive $n$-th derivative tower $D^n f$ (\texttt{fieldDerivN}) obeys $D^m(D^n f)=D^{m+n}f$, annihilates quadratic polynomials at all orders $n\ge 3$ and cubic polynomials at all orders $n\ge 4$ with exact remainder-free Taylor--Newton series $6f_3(x+h)=6f_3(x)+6h\,Df_3(x)+3h(h-1)D^2f_3(x)+h(h-1)(h-2)D^3f_3(x)$, and preserves strictly positive denominators $\prod_{j=0}^n g(x+j)\ge 1$ across all orders $n\in\mathbb{N}$ whenever $g\ge 1$ ($C^\infty$ pole-free smoothness).
\end{enumerate}
\end{theorem}

\section{Native Non-Elementary Basis and Closed-Form Evaluation of Fresnel, Jacobi, Dirichlet, Parseval, and Bessel Integrals}
\label{sec:integrals}

In Liouville's differential algebra, Fresnel integrals $\int\cos(\varphi^2)\,\mathrm{d}\varphi$, elliptic integrals, Dirichlet's sine integral $\mathrm{Si}(\Phi)$, and Bessel integrals $J_0(k)$ admit no closed form in terms of a single circular rotator. On $\mathbb{Z}_{N_{\mathrm{ring}}(d)}$, two native operators close these families:
\begin{enumerate}[leftmargin=*,noitemsep]
\item The two-tick symplectic Verlet difference recurrence
  \begin{equation}
    \label{eq:verlet-parabola}
    \varphi(t+1) - 2\varphi(t) + \varphi(t-1) = \Phi_{\mathrm{accel}}
    \quad\Longrightarrow\quad
    \varphi(t) = \varphi_0 + \omega_0 t + \Phi_{\mathrm{accel}}\,\frac{t(t-1)}{2}
  \end{equation}
  generates Liouville's elementary circular basis at zero acceleration ($\Phi_{\mathrm{accel}}=0$) and the quadratic-phase Fresnel--Airy basis at constant non-zero acceleration ($\Phi_{\mathrm{accel}}\ne 0$).
\item On two coupled rings at the intrinsic cuboctahedron modulus $k^2 = 1/\|e\|_2^2 = 1/(d-1) = 1/2$ (Gauss's lemniscate modulus with complex multiplication in $\mathbb{Z}[i]$), the Jacobi triad $\mathrm{sn}_d(\varphi):=\mathcal{S}_d(\varphi)$, $\mathrm{cn}_d(\varphi):=\mathcal{C}_d(\varphi)$, $2\,\mathrm{dn}_d(\varphi):=\Phi_{\mathrm{half}}(d)-|\mathcal{S}_d(\varphi)|$ satisfies both conservation laws in exact integers.
\end{enumerate}

Define the five canonical ring integrals:
\begin{align}
  \mathcal{C}_{\mathrm{Fresnel}}(d) &:= \sum_{\varphi=0}^{N_{\mathrm{ring}}(d)-1}\mathcal{C}_d(\varphi^2),\qquad
  \mathcal{S}_{\mathrm{Fresnel}}(d) := \sum_{\varphi=0}^{N_{\mathrm{ring}}(d)-1}\mathcal{S}_d(\varphi^2), \label{eq:def-fresnel}\\
  \mathcal{I}_{\mathrm{ell}}(d;a,b) &:= \sum_{\varphi=0}^{N_{\mathrm{ring}}(d)-1}\bigl(a\,|\mathcal{C}_d(\varphi)| + b\,|\mathcal{S}_d(\varphi)|\bigr), \label{eq:def-elliptic}\\
  \mathrm{Si}_d(\Phi) &:= \sum_{\varphi=1}^{\Phi}\mathrm{tdiv}\bigl(\mathcal{S}_d(\varphi),\,\varphi\bigr), \label{eq:def-dirichlet}\\
  \mathcal{I}_{\mathrm{Parseval}}(d) &:= \sum_{\varphi=0}^{N_{\mathrm{ring}}(d)-1}\mathcal{C}_d(\varphi)^2 = \sum_{\varphi=0}^{N_{\mathrm{ring}}(d)-1}\mathcal{S}_d(\varphi)^2, \label{eq:def-parseval}\\
  \mathcal{J}_{0,d}(k) &:= \sum_{\varphi=0}^{N_{\mathrm{ring}}(d)-1}\mathcal{C}_d\bigl(k\,\mathcal{S}_d(\varphi)\bigr). \label{eq:def-bessel}
\end{align}

\begin{theorem}[Lemniscate Jacobi Elliptic Triad Invariants and Parametric Multi-Turn Elliptic Integrals]
\label{thm:main-2}
For all $Q,m,k_{\mathrm{turns}},a,b\in\mathbb{Z}$ and any integration horizon $L\in\mathbb{N}$, the Jacobi elliptic triad on $\mathbb{Z}_{4Q}$ satisfies (\texttt{thm2\_closed\_form\_fresnel\_jacobi\_dirichlet\_parseval\_bessel\_integrals}):
\[
  |\mathrm{sn}_Q(\varphi)| + |\mathrm{cn}_Q(\varphi)| \;=\; m + (Q-m) \;\equiv\; Q,\qquad
  2\,\mathrm{dn}_Q(\varphi) + |\mathrm{sn}_Q(\varphi)| \;=\; (2Q-m) + m \;\equiv\; 2Q,
\]
so factoring the pointwise-constant invariants out of the integral yields $\sum_{\varphi=0}^{L-1}(|\mathrm{sn}_Q|+|\mathrm{cn}_Q|) = L Q$ and $\sum_{\varphi=0}^{L-1}(2\,\mathrm{dn}_Q+|\mathrm{sn}_Q|) = 2 L Q$, while for $d=3$ ($Q_\square(3)=128$) and any linear weights $a,b\in\mathbb{Z}$:
\begin{equation}
  \label{eq:parametric-elliptic}
  k_{\mathrm{turns}}\,\mathcal{I}_{\mathrm{ell}}(3;\,a,b) \;=\; 32,768\,k_{\mathrm{turns}}(a+b),\qquad
  \sum_{\varphi=0}^{127}\bigl(a\,\mathcal{C}_3(\varphi) + b\,\mathcal{S}_3(\varphi)\bigr) \;=\; 8256\,a + 8128\,b.
\end{equation}
\end{theorem}

\begin{theorem}[Closed-Form Fresnel Diffraction, Dirichlet, Parseval, and Bessel Integrals]
\label{thm:main-2b}
For $d=3$ and any amplitude factor $c_{\mathrm{sq}}\in\mathbb{Z}$ and turn count $k_{\mathrm{turns}}\in\mathbb{Z}$ (\texttt{thm2\_closed\_form\_fresnel\_jacobi\_dirichlet\_parseval\_bessel\_integrals}):
\begin{enumerate}[label=\textbf{(\roman*)},leftmargin=*,noitemsep]
\item \textbf{Fresnel diffraction integrals:}
  $\sum_{\varphi=0}^{511} c_{\mathrm{sq}}\mathcal{C}_3(\varphi^2) = \sum_{\varphi=0}^{511} c_{\mathrm{sq}}\mathcal{S}_3(\varphi^2) = 2048\,c_{\mathrm{sq}}$, concentrated at the stationary-phase bisector $\psi = K_P(3) = 64$.
\item \textbf{Dirichlet sine integral and Parseval energy integral:}
  $\mathrm{Si}_3(128) = \mathrm{Si}_3(256) = 128$, and $6\,k_{\mathrm{turns}}\mathcal{I}_{\mathrm{Parseval}}(3) = k_{\mathrm{turns}}N_{\mathrm{ring}}(3)(2Q_\square(3)^2+1) = 16,777,728\,k_{\mathrm{turns}}$ ($\mathcal{I}_{\mathrm{Parseval}}(3) = 2,796,288$).
\item \textbf{Discrete Bessel integral and its exact root:}
  $\mathcal{J}_{0,3}(0) = 65,536$, $\mathcal{J}_{0,3}(1) = 32,768$, and $\mathcal{J}_{0,3}(d_{\mathrm{in}}(\square)) = \mathcal{J}_{0,3}(4) = 0$.
\end{enumerate}
\end{theorem}

\begin{proof}[Proof of Theorems~\ref{thm:main-2}--\ref{thm:main-2b}]
In Theorem~\ref{thm:main-2}, $m+(Q-m)\equiv Q$ and $(2Q-m)+m\equiv 2Q$ are pointwise constants in $m$, so Theorem~\ref{thm:main-1} extracts them out of $\sum_{\varphi=0}^{L-1}$. For Theorem~\ref{thm:main-2b}(i), partitioning $\varphi\in\{0,\dots,511\}$ by $2$-adic valuation cancels all cosets via Euler's identity $\mathcal{E}_3(\psi+256)=-\mathcal{E}_3(\psi)$ except the $32$ elements $\varphi=8(2k+1)$, which satisfy $\varphi^2\equiv 64\pmod{512}$ and yield $32\times 64 = 2048$.
\end{proof}

\section{Emergence of the Continuum Circle, \texorpdfstring{$\pi$}{pi}, and the Riemann Integral from 512-gon Rectification and the Equizonal Partition of the Cuboctahedron}
\label{sec:continuum-emergence}

For any 3D function $F\colon\mathbb{Z}^3\to\mathbb{Z}$, define the natural center-plus-shell integrator on the cuboctahedron ball $B_1(d)=\{0\}\cup O_2(d)$, weighting the origin $(0,0,0)$ by $|V_{O_2}(d)|=12$ equally against the sum over the $12$ vertices of $O_2(d)$:
\begin{equation}
  \label{eq:cuboct-3d-integrator}
  \widehat{\mathcal{I}}_{O_2(d)}(F) \;:=\; |V_{O_2}(d)|\cdot F(0,0,0) \;+\; \sum_{e\in O_2(d)} F(e_x,e_y,e_z).
\end{equation}

\begin{theorem}[Equizonal Partition and Functional 3D Cubature on the Cuboctahedron]
\label{thm:main-3}
For $d=3$, any functions $F,G\colon\mathbb{Z}^3\to\mathbb{Z}$, and any coefficients $c_0,c_x,c_y,c_z,c_{xy},c_{xz},c_{yz}\in\mathbb{Z}$ (\texttt{thm3\_continuum\_integration\_from\_cuboctahedron\_archimedes\_layers}):
\begin{enumerate}[label=\textbf{(\roman*)},leftmargin=*,noitemsep]
\item \textbf{Equizonal partition into Simpson weights $(4,16,4)=4\times(1,4,1)$:}
  slicing $O_2(3)$ along $e_x\in\{-1,0,+1\}$ partitions the $12$ vertices into three equal zones of $d_{\mathrm{in}}(\square)=4$ vertices each; adding the origin weight $|V_{O_2}(3)|=12$ to the equatorial zone $e_x=0$ yields $(4,16,4)$ with total weight $24$.
\item \textbf{Functional additivity, odd-function annihilation, and unified even quadratic/quartic cubature:}
  for all functions $F,G\colon\mathbb{Z}^3\to\mathbb{Z}$ of arbitrary degree,
  \begin{equation}
    \label{eq:functional-odd-annihilation}
    \widehat{\mathcal{I}}_{O_2(3)}(F+G) = \widehat{\mathcal{I}}_{O_2(3)}(F) + \widehat{\mathcal{I}}_{O_2(3)}(G),\quad
    \widehat{\mathcal{I}}_{O_2(3)}\bigl(F(\mathbf{e}) - F(-\mathbf{e})\bigr) = 0,\quad
    \widehat{\mathcal{I}}_{O_2(3)}\bigl(F(e_x,e_y,e_z) - F(-e_x,e_y,e_z)\bigr) = 0,
  \end{equation}
  and on every even quadratic/quartic polynomial ($e_i^4 = e_i^2$ on $O_2(3)$):
  \begin{equation}
    \label{eq:unified-even-cubature}
    \widehat{\mathcal{I}}_{O_2(3)}\!\bigl(c_0 + c_x e_x^2 + c_y e_y^2 + c_z e_z^2 + c_{xy} e_x^2 e_y^2 + c_{xz} e_x^2 e_z^2 + c_{yz} e_y^2 e_z^2\bigr)
    \;=\; 24\,c_0 + 8(c_x+c_y+c_z) + 4(c_{xy}+c_{xz}+c_{yz}),
  \end{equation}
  simultaneously reproducing the exact 3D continuum Riemann integral $3\iiint_{[-1,1]^3} P\,\mathrm{d}V = 24c_0 + 8(c_x+c_y+c_z)$ for every cubic polynomial $P$, the $O(|\mathbf{k}|^2)$ spherical isotropy $\sum_{e\in O_2(3)}(\mathbf{k}\cdot e)^2 = 8|\mathbf{k}|^2$, and the $O(|\mathbf{k}|^4)$ fourth-moment law $\sum_{e\in O_2(3)}(\mathbf{k}\cdot e)^4 = 8\sum k_i^4 + 24\sum_{i<j}k_i^2 k_j^2 = 12|\mathbf{k}|^4 - 4\sum k_i^4$.
\end{enumerate}
\end{theorem}

\begin{theorem}[Wronskian-Cube Rational $\pi(n)$ Brackets, Verlet-Parabola Sine Defect Law, and 3D Euler Norm-Scaling]
\label{thm:main-3b}
For all $Q,p,q,X,Y,x,y\in\mathbb{Z}$ (\texttt{thm3\_continuum\_integration\_from\_cuboctahedron\_archimedes\_layers}):
\begin{enumerate}[label=\textbf{(\roman*)},leftmargin=*,noitemsep]
\item \textbf{Parallelogram radius bound and rational $\pi(n)$ brackets:}
  $2 R_Q(\varphi) = Q^2 + (|\mathcal{C}_Q(\varphi)|-|\mathcal{S}_Q(\varphi)|)^2 \ge Q^2$; on the Weierstrass $n$-gon ($s=N_{\mathrm{ring}}(3)/n$), the inscribed and circumscribed estimators satisfy $\pi_{\mathrm{out}}(n) - \pi_{\mathrm{in}}(n) = \sum_{j=0}^{n/2-1}\frac{W_s(js)^3}{I_s(js)R_3(js)R_3((j+1)s)} > 0$ with strict nested monotonicity $\pi_{\mathrm{in}}(8)<\pi_{\mathrm{in}}(16)<\pi_{\mathrm{in}}(32)<\pi_{\mathrm{out}}(32)<\pi_{\mathrm{out}}(16)<\pi_{\mathrm{out}}(8)$ and Huygens--Snellius value $\pi_3(16)=142/45$.
\item \textbf{Verlet-parabola rational sine defect law and 3D Euler norm-scaling:}
  for the Verlet parabola $P(p,q)=2pq-p^2$, the rational sine numerator $\mathrm{Num}(p,q)=4P(p,q)$ and denominator $\mathrm{Denom}(p,q)=5q^2-P(p,q)$ obey the identity
  \begin{equation}
    \label{eq:verlet-sine-defect}
    \mathrm{Denom}(p,q) - \mathrm{Num}(p,q) \;=\; \bigl(5q^2 - (2pq-p^2)\bigr) - 4(2pq-p^2) \;=\; 5(q-p)^2 \;\ge\; 0,
  \end{equation}
  while every planar Weierstrass rotation across the $3$ cuboctahedron coordinate planes scales squared Euclidean norms identically via $(Xx-Yy)^2 + (Yx+Xy)^2 = (X^2+Y^2)(x^2+y^2)$.
\end{enumerate}
\end{theorem}

\section{Formal Verification in Lean~4}
\label{sec:lean4}

The formalization of Theorems~\ref{thm:main-1}--\ref{thm:main-3b} (\texttt{thm1\_constant\_wronskian\_weierstrass\_lift\_and\_newton\_leibniz}, \texttt{thm2\_closed\_form\_fresnel\_jacobi\_dirichlet\_parseval\_bessel\_integrals}, \texttt{thm3\_continuum\_integration\_from\_cuboctahedron\_archimedes\_layers}, \texttt{thm34\_rotator\_projective\_field\_invertibility\_zero\_divisor\_freedom\_and\_q\_i\_sqrt2\_fields}, \texttt{thm35\_general\_field\_derivative\_operator\_quotient\_wronskian\_and\_variational\_euler\_lagrange}, and \texttt{thm38\_c\_infinity\_higher\_derivatives\_taylor\_newton\_expansion\_and\_pole\_free\_quotients}) resides in \texttt{RotatorQuadrature.lean} in \url{https://github.com/AI-Guiders/scientific-papers} (\texttt{0 sorryAx}, \texttt{0 Classical.choice}).

\begin{thebibliography}{9}
\bibitem{Boole1860}
G.~Boole, \emph{A Treatise on the Calculus of Finite Differences}, Macmillan and Co., Cambridge--London, 1860.
\bibitem{Malcev1973}
A.\,I.~Mal'cev, \emph{Algebraic Systems}, Die Grundlehren der mathematischen Wissenschaften, vol.~192, Springer-Verlag, Berlin--Heidelberg--New York, 1973 (Russian original: \emph{Алгебраические системы}, Nauka, Moscow, 1970).
\bibitem{Sobolev1974}
S.\,L.~Sobolev, \emph{Introduction to the Theory of Cubature Formulae}, Nauka, Moscow, 1974 (English transl.: \emph{Cubature Formulas and Modern Analysis}, Gordon and Breach, 1992).
\bibitem{Krause1986}
E.\,F.~Krause, \emph{Taxicab Geometry: An Adventure in Non-Euclidean Geometry}, Dover Publications, New York, 1986.
\bibitem{Thompson2011}
K.\,P.~Thompson, ``The nature of length, area, and volume in taxicab geometry,'' \emph{arXiv preprint} \texttt{arXiv:1101.2922} [math.MG], 2011.
\bibitem{Thompson2012}
K.\,P.~Thompson, ``Taxicab calculus: Trig derivatives,'' \emph{arXiv preprint} \texttt{arXiv:1205.3814v2} [math.MG], 2012.
\bibitem{Stroud1971}
A.\,H.~Stroud, \emph{Approximate Calculation of Multiple Integrals}, Prentice-Hall, Englewood Cliffs, NJ, 1971.
\bibitem{Moura2021}
L.~de~Moura and S.~Ullrich, ``The Lean~4 theorem prover and programming language,'' in \emph{CADE~28}, LNCS vol.~12699, Springer, 2021, pp.~625--635.
\end{thebibliography}

\end{document}
